\documentclass{article}
\usepackage[T1]{fontenc}
\usepackage{lmodern}
\usepackage{graphicx}
\usepackage{amsmath,amssymb}
\usepackage[a4paper,margin=2cm]{geometry}
\usepackage[absolute,overlay]{textpos}
\setlength{\TPHorizModule}{1mm}
\setlength{\TPVertModule}{1mm}
\usepackage[indonesian]{babel}
\usepackage{hyperref}
\setlength{\emergencystretch}{2em}
% unit: Halaman 69

\begin{document}

% === Halaman 69 (llm) ===

\begin{itemize}
  \item Sifat \textit{reversible} ini membuat perancang cipher blok tidak perlu membuat algoritma baru untuk mendekripsi cipherteks menjadi plainteks.

  Contoh: \[ L_{i-1} \oplus f(R_{i-1}, K_i) \oplus f(R_{i-1}, K_i) = L_{i-1} \]

  \item Sifat \textit{reversible} tidak bergantung pada fungsi $f$ sehingga fungsi $f$ dapat dibuat serumit mungkin.
\end{itemize}

\end{document}
